\section{Research Protocol}
\label{sec:protocol}

\subsection{Data source and admission}
The authorized owner copy of the dataset is at \texttt{/data/yingda/tidetail-data/prime-formal-2015-2024},
containing 350 files totaling 451~MB. Six coastal regions are included, with
87,672 hourly UTC rows each spanning 2015-01-01 to 2024-12-31. Region admission
status is recorded in \texttt{region\_availability.json}:

\begin{table}[htbp]
\centering
\caption{Region formal training authorization status}
\begin{tabular}{l c}
\toprule
Region & Training authorized \\
\midrule
New York Bight & Yes \\
Thames / North Sea & Yes \\
Hong Kong Pearl River & \textbf{No} (license/datum not signed) \\
Southeast Australia & Yes \\
US Gulf Coast & Yes \\
Charleston Atlantic & Yes \\
\bottomrule
\end{tabular}
\label{tab:regions}
\end{table}

The Hong Kong region is excluded from all formal experiments until the data owner
confirms license, vertical datum, and the HK event \texttt{pearl\_river\_hong\_kong\_event\_01149},
which spans the validation/test boundary at index 78904. The migration manifest
(\texttt{MIGRATION\_MANIFEST.json}) exists one level above the release directory;
file SHA256 and size checks confirm all inspected files match both release and
migration inventory entries.

\subsection{Target definition}
The release manifest \texttt{database\_manifest.json} defines:
\begin{itemize}
    \item \textbf{Total water level}: observed water level above local chart datum.
    \item \textbf{Tide}: deterministic FES2022 tidal constituent forecast (no future
    uncertainty in this archival release).
    \item \textbf{Residual}: $\text{water\_level\_m} = \text{observed\_total} - \text{FES2022\_tide}$.
\end{itemize}
The pilot diffusion model targets \texttt{water\_level\_m} (residual), in metres,
with a fixed 72-hour history $\rightarrow$ 24-hour forecast window for engineering
feasibility. The formal release contract specifies 96-hour history $\rightarrow$
72-hour forecast; all comparisons between methods should use one consistent
contract for scientific validity.

\subsection{Leakage-safe splitting}
Our protocol enforces three hard gates before any window creation or normalization:

\begin{enumerate}
    \item \textbf{Event-group disjointness}: the same independent storm event may
    not appear in two partitions. Events are represented as \texttt{EventGroup(event\_id,
    start\_idx, end\_idx)} intervals. The implementation in \texttt{src/tidetail/splits.py}
    validates non-overlapping intervals and unique IDs before splitting.

    \item \textbf{Chronological ordering}: train row block must end before
    validation row block, which must end before test row block. The split is
    8:1:1 by chronological event-group order, subject to \texttt{min\_events\_per\_split}.

    \item \textbf{Full-support embargo}: for forecast origin time $o$, lookback $L$,
    and horizon $H$, the full temporal support of one forecast window is
    $[o-L, o+H)$. Two partitions must not have overlapping full-support intervals.
    The repaired \texttt{embargo\_filter} compares every index in an event group
    to every protected index, replacing the earlier midpoint-only comparison that
    incorrectly accepted dangerous events (see Appendix~\ref{sec:audit}).
    Additionally, \texttt{embargo\_full\_support\_filter} enforces
    $o \geq \text{split\_block\_start} + L$ and $o \leq \text{split\_block\_end} - H$
    for each forecast origin.
\end{enumerate}

All scaler fitting, quantile threshold estimation, and interpolation fitting are
performed on \emph{train rows only} and then applied to all partitions without
re-fitting. Threshold-run events crossing the train boundary are removed in full
before window creation, not partial-interval trimmed.

\subsection{Formal release contract}
The database manifest specifies:
\begin{itemize}
    \item \textbf{Input}: 96 hours of historical observations.
    \item \textbf{Horizon}: 72 hours of forecast.
    \item \textbf{Split}: chronological 8:1:1 with strict window containment.
    \item \textbf{Tide}: existing local FES2022 deterministic archive.
    \item \textbf{Forcing}: ERA5 hindcast only. No operational 72-hour forecast
    claim without archived issue/valid/lead NWP forecast vintages.
\end{itemize}
The ERA5 forcing data available in this release is historical hindcast data, not
forecast data. Using future observed ERA5 values as conditioning information in
a forecast experiment would constitute leakage unless explicitly labelled as an
oracle experiment. Our gauge-only diffusion model does not consume atmospheric
forcing; this limitation is discussed in Section~\ref{sec:discussion}.
